% TEMPLATE-MILC-v3.tex - plantilla maestra para todas las guias MILC v3
% Ver scripts/generadores/build-guia-milc.py para la lista completa de
% claves esperadas. El script valida que no queden placeholders sin
% reemplazar antes de escribir el archivo final.

\documentclass[11pt,letterpaper]{article}
\usepackage[margin=1.75cm]{geometry}
\usepackage[table]{xcolor}
\usepackage{fontspec}
\usepackage[spanish]{babel}
\usepackage{tikz}
\usepackage[most]{tcolorbox}
\usepackage{tabularx}
\usepackage{array}
\usepackage{fancyhdr}
\usepackage{titlesec}
\usepackage{ragged2e}
\usepackage{enumitem}
\usepackage{microtype}

\setmainfont{Helvetica Neue}
\setsansfont{Helvetica Neue}
\setlength{\parindent}{0pt}
\setlength{\parskip}{6pt}
\hyphenpenalty=200
\exhyphenpenalty=200
\pretolerance=400
\tolerance=2000
\setlength{\emergencystretch}{1em}
\renewcommand{\arraystretch}{1.25}
\setlist{leftmargin=5mm,itemsep=1mm,topsep=1mm}
\newcolumntype{Y}{>{\RaggedRight\arraybackslash}X}

\definecolor{milcMagenta}{HTML}{E6007E}
\definecolor{milcVino}{HTML}{5A0038}
\definecolor{milcTurquesa}{HTML}{0093A5}
\definecolor{milcVerde}{HTML}{58A923}
\definecolor{milcMostaza}{HTML}{E5B400}
\definecolor{milcOcre}{HTML}{B86600}
\definecolor{milcNegro}{HTML}{241816}
\definecolor{milcGris}{HTML}{F7F7F4}
\definecolor{milcLinea}{HTML}{E8E2DD}
\definecolor{milcGradePrimary}{HTML}{0066FF}
\definecolor{milcGradeDark}{HTML}{003D99}
\definecolor{milcGradeSoft}{HTML}{D6E8FF}
\definecolor{milcGradeAccent}{HTML}{CFE7FF}
\definecolor{milcGradeText}{HTML}{FFFFFF}
\color{milcNegro}

\pagestyle{fancy}
\fancyhf{}
\lhead{\small Tecnología e Informática · Grado 6}
\rhead{\small Guía 13 · MILC v3}
\cfoot{\small\thepage}
\renewcommand{\headrulewidth}{0pt}

\newtcolorbox{titlebox}[1]{
  enhanced,colback=#1,colframe=#1,arc=7mm,boxrule=0pt,
  left=7mm,right=7mm,top=3mm,bottom=3mm,
  before skip=8pt,after skip=8pt,
  fontupper=\sffamily\bfseries\Large\color{white}
}

\newtcolorbox{softbox}[3]{
  enhanced,colback=#2,colframe=#1,borderline west={4pt}{0pt}{#1},
  arc=2mm,boxrule=.4pt,left=4mm,right=4mm,top=3mm,bottom=3mm,
  before skip=6pt,after skip=8pt,title={#3},coltitle=white,
  colbacktitle=#1,fonttitle=\sffamily\bfseries,fontupper=\RaggedRight,
  attach boxed title to top left={xshift=3mm,yshift=-2mm},
  boxed title style={arc=2mm, boxrule=0pt}
}

\newtcolorbox{drawbox}[2]{
  enhanced,colback=white,colframe=#1,arc=4mm,boxrule=1pt,
  left=5mm,right=5mm,top=4mm,bottom=4mm,before skip=8pt,after skip=8pt,
  title={#2},coltitle=white,colbacktitle=#1,fonttitle=\sffamily\bfseries,
  fontupper=\RaggedRight,
  attach boxed title to top left={xshift=4mm,yshift=-2mm},
  boxed title style={arc=3mm, boxrule=0pt}
}

\newtcolorbox{infoband}[1]{
  enhanced,colback=#1!8,colframe=#1!50,arc=2mm,boxrule=.5pt,
  left=4mm,right=4mm,top=2mm,bottom=2mm,before skip=4pt,after skip=4pt,
  fontupper=\small\RaggedRight
}

% Puente narrativo entre secciones (contrato editorial Conéctate).
% Da continuidad: "Acabas de X. Ahora vas a Y porque Z."
% Visualmente: banda izquierda en color del grado, texto en itálica pequeña.
\newtcolorbox{puentebox}{
  enhanced,colback=milcGradeSoft,colframe=milcGradeDark,
  borderline west={3pt}{0pt}{milcGradeDark},
  arc=0mm,boxrule=0pt,left=5mm,right=4mm,top=2.5mm,bottom=2.5mm,
  before skip=4pt,after skip=8pt,
  fontupper=\itshape\small\color{milcNegro}\RaggedRight
}
\newcommand{\puente}[1]{%
  \begin{puentebox}%
    \textcolor{milcGradeDark}{\textbf{\textrightarrow}}\hspace{2mm}#1%
  \end{puentebox}%
}

\newcommand{\linead}[1]{\noindent\color{milcLinea}\rule{#1}{0.5pt}\color{milcNegro}}
\newcommand{\respuesta}[1]{\par\vspace{2mm}\foreach \n in {1,...,#1}{\noindent\color{milcLinea}\rule{\linewidth}{0.5pt}\par\vspace{5mm}}\color{milcNegro}}
\newcommand{\checkbox}{\textcolor{milcGradeDark}{\fbox{\phantom{x}}}\hspace{5pt}}

\newcommand{\phasecard}[4]{
\begin{tcolorbox}[enhanced,colback=#3,colframe=#2,arc=3mm,boxrule=.5pt,height=3.05cm,valign=center,halign=center,left=1.5mm,right=1.5mm,top=2mm,bottom=2mm]
{\sffamily\bfseries\fontsize{22}{24}\selectfont\textcolor{#2}{#1}}\par
{\sffamily\bfseries\small #4}
\end{tcolorbox}
}

\begin{document}
\thispagestyle{empty}
\begin{tikzpicture}[remember picture,overlay]
  \fill[white] (current page.south west) rectangle (current page.north east);
  \path[fill=milcGradePrimary,rounded corners=16mm]
    ([xshift=.55cm,yshift=-.55cm]current page.north west) rectangle
    ([xshift=-.55cm,yshift=3.65cm]current page.south east);

  \node[anchor=north west,text=milcGradeText] at ([xshift=2.0cm,yshift=-1.85cm]current page.north west)
    {\sffamily\bfseries\fontsize{14}{16}\selectfont GRADO};
  \node[anchor=north west,text=milcGradeText,opacity=.82] at ([xshift=2.0cm,yshift=-2.75cm]current page.north west)
    {\sffamily\bfseries\fontsize{9}{11}\selectfont SERIE GUÍAS · TECNOLOGÍA E INFORMÁTICA};

  \begin{scope}[shift={([xshift=-3.05cm,yshift=-2.55cm]current page.north east)}]
    \fill[white,opacity=.80] (-.62,-.52) rectangle (.62,.52);
    \draw[milcGradeText,line width=1pt] (-.62,-.52) rectangle (.62,.52);
    \fill[milcGradeDark] (-.42,-.38) rectangle (-.22,.06);
    \fill[milcGradeAccent] (-.10,-.38) rectangle (.10,.24);
    \fill[milcGradePrimary!60!black] (.22,-.38) rectangle (.42,.38);
  \end{scope}

  \node[anchor=west,text=milcGradeText] at ([xshift=1.9cm,yshift=-12.85cm]current page.north west)
    {\sffamily\bfseries\fontsize{124}{124}\selectfont 6};
  \draw[milcGradeText,line width=1.7pt]
    ([xshift=4.52cm,yshift=-11.68cm]current page.north west) circle (.13cm);

  \node[anchor=west,text=milcGradeText,text width=8.7cm,align=left] at ([xshift=10.35cm,yshift=-11.6cm]current page.north west)
    {
     {\sffamily\bfseries\fontsize{19}{22}\selectfont El gabinete\\por dentro ---\\conocer las\\piezas como\\el relojero}\\[4mm]
     {\sffamily\bfseries\fontsize{23}{26}\selectfont Guía 13-6-TIC}\\[2mm]
     {\sffamily\fontsize{13}{16}\selectfont Período 2 · Hardware y software}\\[5mm]
     {\sffamily\bfseries\fontsize{11}{13}\selectfont Institución Educativa Sor María Juliana}\\[1mm]
     {\sffamily\fontsize{10}{12}\selectfont Autor: Dr. Álvaro Cárdenas Orozco}
    };

  \path[fill=milcGradeSoft,rounded corners=7mm]
    ([xshift=1.35cm,yshift=.85cm]current page.south west) rectangle
    ([xshift=-1.35cm,yshift=4.25cm]current page.south east);
  \draw[milcGradeDark,line width=.65pt,rounded corners=7mm]
    ([xshift=1.35cm,yshift=.85cm]current page.south west) rectangle
    ([xshift=-1.35cm,yshift=4.25cm]current page.south east);
  \node[anchor=center,text=milcNegro,text width=17.0cm,align=center] at ([yshift=2.55cm]current page.south)
    {\sffamily\fontsize{10}{12}\selectfont Metodología MILC · Apertura ancestral · Pensamiento computacional · Triángulo Dussel-Estoicismo-Floridi\\
    \textcolor{milcGradeDark}{\bfseries Producto:} Un \textbf{plano del interior del gabinete del computador} dibujado en una hoja completa del cuaderno. El plano incluye \textbf{6 piezas ubicadas correctamente} (tarjeta madre, CPU con disipador, RAM, disco duro, fuente de poder, tarjeta gráfica), cada una con \textbf{una etiqueta} y \textbf{una frase de qué hace}. Más \textbf{una tabla comparativa} de CPU vs RAM vs disco duro: en qué se parecen, en qué se diferencian, cuál es ``cerebro'', cuál es ``memoria corta'' y cuál es ``memoria larga''.};
\end{tikzpicture}
\mbox{}
\newpage

\section*{Apertura: El gabinete por dentro --- conocer las piezas como las conoce el relojero}

\begin{softbox}{milcGradeDark}{milcGradeSoft}{Saber ancestral}
\textbf{Cuando el relojero levantaba la tapa del reloj, mostraba un universo pequeño.} En la calle 14 de Cartago, don Lucho tenía un cristal de aumento (\emph{lupa}) y una pinza diminuta. Cuando llegaba un cliente con un reloj atrasado, abría la tapa frente a él, mostraba los engranajes y \textbf{explicaba cada pieza con paciencia}: \emph{``Mire, este es el balancín, hace el tic-tac. Esta rueda grande mueve las manecillas. Este resortico guarda la cuerda. Este eje es el corazón del mecanismo, sin él nada se mueve.''}. La gente miraba hipnotizada. Esa pequeña clase del relojero \textbf{convertía un misterio en algo familiar}. Cuando lo entendían, ya no le tenían miedo al reloj. Hoy abrimos juntos el gabinete del computador. Lo que parece complicado se va a sentir cercano cuando sepas el nombre de cada pieza y para qué sirve.
\end{softbox}

\begin{softbox}{milcTurquesa}{milcGris}{Saber contemporáneo}
El \textbf{gabinete} (o \emph{torre}, o \emph{case}) es la caja del computador. En los portátiles está \emph{adentro} (no se ve), pero en los de escritorio se puede abrir para ver las piezas. Dentro hay \textbf{6 piezas principales} que toda guía técnica menciona: (1) \textbf{Tarjeta madre} (motherboard) --- la placa grande donde van las demás. (2) \textbf{CPU} (procesador) --- el cerebro. (3) \textbf{RAM} (memoria) --- la mesa de trabajo. (4) \textbf{Disco duro} (HDD o SSD) --- el archivo. (5) \textbf{Tarjeta gráfica} (GPU) --- para los juegos y videos. (6) \textbf{Fuente de poder} --- alimenta de electricidad a todo. \textbf{Reconocerlas y saber dónde van te abre la puerta} a entender cuando hablen de \emph{``actualizar la RAM''} o \emph{``cambiar el disco''}. Es vocabulario que vas a usar el resto de tu vida.
\end{softbox}

\begin{softbox}{milcMostaza}{milcGris}{Pregunta-puente}
\textit{Cuando una persona dice \emph{``mi computador está lento, le voy a meter más RAM''}, ¿\textbf{sabes qué quiere decir} esa frase? ¿Qué es la RAM y por qué meter más ayuda?}
\end{softbox}

\begin{softbox}{milcVerde}{milcGris}{Saber y saber hacer}
Al terminar la clase: (1) podrás \textbf{identificar} las 6 piezas internas del computador; (2) sabrás \textbf{explicar} qué hace cada una en lenguaje natural; (3) podrás \textbf{distinguir} CPU, RAM y disco duro con sus diferencias clave; (4) habrás \textbf{creado} un plano del interior del gabinete bien etiquetado.
\end{softbox}

\small
\begin{tabularx}{\textwidth}{>{\bfseries}p{3.0cm}Y}
\rowcolor{milcGradePrimary!12}Autor & Dr. Álvaro Cárdenas Orozco\\
DBA & Identifica componentes y sistemas tecnológicos en su entorno (MEN, T\&I 6°)\\
Referentes & Saber ancestral del relojero · Sistemas como cadena de oficios · Diagnóstico paso a paso\\
Producto final & Un \textbf{plano del interior del gabinete del computador} dibujado en una hoja completa del cuaderno. El plano incluye \textbf{6 piezas ubicadas correctamente} (tarjeta madre, CPU con disipador, RAM, disco duro, fuente de poder, tarjeta gráfica), cada una con \textbf{una etiqueta} y \textbf{una frase de qué hace}. Más \textbf{una tabla comparativa} de CPU vs RAM vs disco duro: en qué se parecen, en qué se diferencian, cuál es ``cerebro'', cuál es ``memoria corta'' y cuál es ``memoria larga''.\\
\end{tabularx}
\normalsize

\puente{Hoy haces 4 cosas: identificas 6 piezas con sus nombres, aprendes qué hace cada una, distingues CPU/RAM/disco, y dibujas el plano del gabinete.}

\begin{titlebox}{milcGradeDark}
Ruta MILC: así vamos a aprender
\end{titlebox}
\begin{tabularx}{\textwidth}{>{\centering\arraybackslash}X>{\centering\arraybackslash}X>{\centering\arraybackslash}X>{\centering\arraybackslash}X}
\phasecard{1}{milcVerde}{milcGris}{Escuta\\[-1mm]\small Observo, escucho y pregunto.} &
\phasecard{2}{milcTurquesa}{milcGris}{Sistematización\\[-1mm]\small Pensamiento computacional.} &
\phasecard{3}{milcMagenta}{milcGris}{Praxis\\[-1mm]\small Construyo y aplico.} &
\phasecard{4}{milcVino}{milcGris}{Evaluación\\[-1mm]\small Reflexión liberadora.}
\end{tabularx}


\puente{Empezamos jugando a adivinar nombres con 6 imágenes mentales.}

\begin{titlebox}{milcVerde}
Fase 1 · Escuta
\end{titlebox}

\begin{softbox}{milcVerde}{milcGris}{Escena para escuchar}
\textbf{Actividad 1 · IDENTIFICA --- 6 imágenes mentales} (10 min · individual). En el cuaderno, escribe la \textbf{lista de 6 piezas} que se mencionaron en la introducción: tarjeta madre, CPU, RAM, disco duro, tarjeta gráfica, fuente de poder. Al lado de cada una, \textbf{en una sola línea}, escribe \emph{qué crees que hace} con tus palabras --- sin haber leído nada todavía. Es tu intuición. \textbf{Tip:} si no tienes idea, escribe lo que \emph{te suena} por el nombre. Después comparas con lo que vas a aprender. No se vale ver la siguiente actividad.
\end{softbox}

\checkbox Escribí las 6 piezas en orden.\\
\checkbox Al lado de cada una, escribí mi intuición de qué hace.\\
\checkbox No miré la siguiente actividad ni internet.

\begin{infoband}{milcVerde}
\textbf{Lo que va al cuaderno (Actividad 1 · IDENTIFICA).} Título: «Actividad 1 --- Mis 6 intuiciones». \textbf{Formato:} lista numerada del 1 al 6. \textbf{Extensión:} 6 líneas. \textbf{Sabes que terminaste cuando} tienes las 6 piezas con tu propia intuición al lado.
\end{infoband}

\puente{Después del juego, conoces cada pieza con su función y un dato fácil de recordar.}

\begin{titlebox}{milcTurquesa}
Fase 2 · Sistematización con pensamiento computacional
\end{titlebox}

\begin{softbox}{milcTurquesa}{milcGris}{Cuatro pilares aplicados al tema}
Ahora compara tus 6 intuiciones con la realidad. Vas a ver que algunas las acertaste (los nombres dicen mucho) y otras te sorprenden. Lo importante es \textbf{quedarte con el dato fácil de recordar de cada pieza}.
\end{softbox}

\small
\begin{tabularx}{\textwidth}{>{\bfseries\RaggedRight\arraybackslash}p{3.6cm}Y}
\rowcolor{milcTurquesa!90}\color{white}Pilar & \color{white}\bfseries Aplicación al tema\\
1. Descomposición & \textbf{Pieza 1 --- Tarjeta madre (motherboard).} Es la \textbf{placa grande verde} (a veces negra) donde se conectan TODAS las demás piezas. Sin tarjeta madre no hay computador, porque es como la \emph{calle principal del barrio} que conecta todas las casas. \emph{Dato fácil:} se llama ``madre'' porque sostiene y conecta a todas las otras piezas, como una madre con sus hijos. \textbf{Pieza 2 --- CPU (procesador).} Es el \textbf{cerebro} del computador. Un chip cuadrado pequeño (3-5 cm), pegado a la tarjeta madre. Lleva un \emph{disipador} (un ventilador encima) porque trabaja tanto que se calienta. \emph{Dato fácil:} CPU = ``Central Processing Unit'' = ``el que piensa''.\\
2. Reconocimiento de patrones & \textbf{Pieza 3 --- RAM (memoria de trabajo).} Son \textbf{tarjetas alargadas y delgadas}, como reglitas paradas, conectadas a la tarjeta madre. \emph{Sirven como la mesa de trabajo}: la CPU pone ahí lo que está usando AHORA. Más RAM = más mesa = más cosas puedes hacer al mismo tiempo. \emph{Dato fácil:} se borra al apagar el computador. \textbf{Pieza 4 --- Disco duro (HDD o SSD).} Es donde se \textbf{guardan tus archivos para siempre}. HDD = \emph{disco mecánico} (lento, barato, suena cuando gira). SSD = \emph{disco de estado sólido} (rápido, caro, silencioso). \emph{Dato fácil:} aunque apagues el equipo, los archivos quedan acá.\\
3. Abstracción & \textbf{Pieza 5 --- Tarjeta gráfica (GPU).} Es \textbf{otra placa con su propio chip}, especializada en gráficos y videos. Los juegos pesados (Roblox, Minecraft, Fortnite) la necesitan. En computadores básicos no hay una separada; está incluida dentro de la CPU. \emph{Dato fácil:} GPU = ``Graphics Processing Unit'' = ``cerebro para los gráficos''. \textbf{Pieza 6 --- Fuente de poder (PSU).} Es la \textbf{caja metálica grande con cables saliendo}, generalmente en una esquina del gabinete. \emph{Convierte la electricidad de la pared (110V) en electricidad para las piezas}. \emph{Dato fácil:} si la fuente falla, NADA enciende. Es como el corazón que bombea sangre.\\
4. Algoritmo conceptual & \textbf{Comparación rápida CPU vs RAM vs Disco duro.} (1) \textbf{CPU} es \emph{cerebro}: piensa, calcula, decide. (2) \textbf{RAM} es \emph{mesa de trabajo}: guarda lo de AHORA, se borra al apagar. (3) \textbf{Disco duro} es \emph{archivo del cuaderno}: guarda los archivos para siempre, no se borra. \textbf{Analogía completa:} estudiar es como usar el computador. Tu \emph{cerebro} (CPU) piensa. Pones los apuntes que necesitas en tu \emph{mesa} (RAM). Tus \emph{cuadernos archivados} en el cajón (disco duro). Cuando terminas el día, recoges la mesa (la RAM se borra) pero los cuadernos quedan (el disco duro guarda).\\
\end{tabularx}
\normalsize

\begin{drawbox}{milcTurquesa}{Las 6 piezas del gabinete + 3 piezas estrella}
\textbf{1. Tarjeta madre} --- la calle principal que conecta todo. \\[1mm]
\textbf{2. CPU} --- el cerebro. \\[1mm]
\textbf{3. RAM} --- la mesa de trabajo (memoria corta). \\[1mm]
\textbf{4. Disco duro} --- el archivo (memoria larga). \\[1mm]
\textbf{5. Tarjeta gráfica} --- cerebro de los juegos. \\[1mm]
\textbf{6. Fuente de poder} --- el corazón eléctrico. \\[1mm]
\textbf{Las 3 estrella:} CPU (cerebro), RAM (mesa), disco (archivo).
\end{drawbox}

\begin{softbox}{milcMostaza}{milcGris}{Errores comunes y su corrección}
\textbf{Confusiones que se repiten cuando ves piezas por primera vez.} (1) \emph{``RAM y disco duro son lo mismo''} --- Falso. RAM se borra al apagar; disco duro queda. (2) \emph{``Más RAM hace el computador más rápido en todo''} --- Falso. Más RAM ayuda a tener \emph{muchas cosas abiertas a la vez}; para velocidad pura, importa más la CPU. (3) \emph{``Si cambio el disco, pierdo mis archivos''} --- Cierto, los archivos están en el disco. Por eso siempre hay que hacer \emph{copia de seguridad} antes de cambiarlo. (4) \emph{``Los portátiles no tienen estas piezas''} --- Falso. Las tienen todas, solo que más pequeñas y dentro de la carcasa.
\end{softbox}

\begin{infoband}{milcTurquesa}
\textbf{Lo que va al cuaderno (Actividad 2 · EXPLICA).} Título: «Actividad 2 --- Las 6 piezas con su dato fácil». \textbf{Formato:} tabla de 6 filas, 3 columnas. Columna 1: nombre de la pieza. Columna 2: qué hace (1 frase). Columna 3: dato fácil para recordar. \textbf{Extensión:} 6 filas. \textbf{Sabes que terminaste cuando} podrías recitar de memoria las 6 piezas y sus datos fáciles.
\end{infoband}

\puente{Con las 6 piezas claras, dibujas el plano del gabinete y haces la tabla comparativa.}

\begin{titlebox}{milcMagenta}
Fase 3 · Praxis con pensamiento computacional
\end{titlebox}

\begin{softbox}{milcMagenta}{milcGris}{Construyendo el producto con los cuatro pilares}
\textbf{Actividad 3 · CREA --- Tu plano del gabinete + tabla comparativa} (30 min · individual). Vas a hacer 2 cosas en una hoja completa del cuaderno. \textbf{Arriba (3/4 de la hoja):} el plano del interior del gabinete con las 6 piezas dibujadas y etiquetadas. \textbf{Abajo (1/4 de la hoja):} la tabla comparativa de CPU vs RAM vs disco duro.
\end{softbox}

\small
\begin{tabularx}{\textwidth}{>{\bfseries\RaggedRight\arraybackslash}p{3.6cm}Y}
\rowcolor{milcMagenta!90}\color{white}Pilar & \color{white}\bfseries Aplicación al producto\\
Descomposición & \textbf{Paso 1 --- Título y marco.} En la parte superior escribe: \textbf{``Plano del gabinete''} y debajo: \emph{``[Tu nombre], grado 6[tu letra], año 2026''}. Subraya el título. Dibuja un \textbf{rectángulo grande vertical} que ocupe los 3/4 superiores de la hoja --- ese es el gabinete visto desde un lado abierto.\\
Patrones & \textbf{Paso 2 --- Ubica la tarjeta madre.} En el centro del rectángulo dibuja una \textbf{placa rectangular grande} (de color verde si tienes lápices de color, si no, solo el contorno). Etiqueta: \emph{``1. Tarjeta madre''}. Esta es la base sobre la que van las demás.\\
Abstracción & \textbf{Paso 3 --- CPU + disipador.} Encima de la tarjeta madre dibuja un \textbf{cuadrado pequeño con un cuadrado más grande encima} (el cuadrado de encima es el disipador con su ventilador). Etiqueta: \emph{``2. CPU --- cerebro''}. La CPU está casi en el centro de la tarjeta madre.\\
Algoritmo de armado & \textbf{Paso 4 --- RAM.} A la derecha de la CPU dibuja \textbf{2 reglitas paradas} (2 ó 4 módulos de RAM). Etiqueta: \emph{``3. RAM --- mesa de trabajo''}. La RAM siempre va a un lado de la CPU, conectadas por ranuras de la tarjeta madre.\\
\end{tabularx}
\normalsize

\begin{drawbox}{milcMagenta}{Antes de mostrar tu plano}
\checkbox El plano ocupa los 3/4 superiores de la hoja. \\[1mm]
\checkbox Están las 6 piezas dibujadas y etiquetadas. \\[1mm]
\checkbox La tarjeta madre está en el centro y las demás encima o conectadas. \\[1mm]
\checkbox La tabla comparativa tiene 3 columnas y 4 filas. \\[1mm]
\checkbox La tabla deja claro cuál es cerebro, cuál memoria corta y cuál memoria larga. \\[1mm]
\checkbox Está firmado con nombre y fecha.
\end{drawbox}

\begin{softbox}{milcMagenta}{milcGris}{Plantilla de trabajo}
\textbf{Estructura.} \textit{Arriba:} título + nombre. \textit{Centro (rectángulo grande):} tarjeta madre + CPU con disipador + RAM + disco + tarjeta gráfica + fuente, las 6 etiquetadas. \textit{Abajo:} tabla CPU/RAM/disco con 4 filas (qué hace, memoria, ¿se borra al apagar?, cómo recordarlo). \textit{Esquina:} firma.
\end{softbox}

\begin{infoband}{milcMagenta}
\textbf{Lo que va al cuaderno (Actividad 3 · CREA).} Título: «Actividad 3 --- Mi plano del gabinete». \textbf{Formato:} plano dibujado + tabla comparativa. \textbf{Extensión:} 1 hoja entera. \textbf{Sabes que terminaste cuando} podrías mostrarle el plano a un primo y él entendería las 6 piezas y cuál es cerebro, cuál mesa y cuál archivo.
\end{infoband}

\puente{El producto es el plano firmado + la tabla CPU vs RAM vs disco.}

\begin{titlebox}{milcMostaza}
Producto de la sesión
\end{titlebox}
\textbf{Plano del gabinete con 6 piezas + tabla CPU vs RAM vs disco · firmado}.
\begin{softbox}{milcMostaza}{milcGris}{Criterios de calidad}
(1) El plano contiene las 6 piezas dibujadas con etiqueta. (2) La tarjeta madre está como base y las demás conectadas a ella. (3) Cada pieza tiene una frase breve de qué hace. (4) La tabla comparativa distingue claramente cerebro, memoria corta y memoria larga. (5) Está firmado con nombre y fecha.
\end{softbox}

\puente{Antes de salir, verificas que tu plano tiene las 6 piezas y que la tabla aclara cuál es cerebro, memoria corta y memoria larga.}

\begin{titlebox}{milcVino}
Fase 4 · Evaluación liberadora — cinco dimensiones
\end{titlebox}

\begin{softbox}{milcVino}{milcGris}{Sentido de la evaluación}
Evaluar no es solo revisar una nota. Es reconocer qué aprendí, qué cambió en mí, qué emociones manejé, cómo me relaciono con la ciudad y la comunidad, y qué le dejaré a quien viene detrás.
\end{softbox}

\textbf{Mi reflexión liberadora en cinco dimensiones.} Completa cada frase iniciada:

\small
\begin{tabularx}{\textwidth}{>{\bfseries\RaggedRight\arraybackslash}p{3.4cm}Y>{\RaggedRight\arraybackslash}p{4.6cm}}
\rowcolor{milcVino}\color{white}Dimensión & \color{white}\bfseries Mi respuesta (frame starter) & \color{white}\bfseries Algo que descubrí\\
1. Personal & Lo que más cambió en mí fue \linead{6.5cm} & \linead{4.2cm}\\
2. Emocional & La emoción más fuerte fue \linead{2.5cm} y la regulé \linead{3cm} & \linead{4.2cm}\\
3. Ciudadana & En democracia esto importa porque \linead{6cm} & \linead{4.2cm}\\
4. Local (Cartago) & En mi barrio sería útil para \linead{6.5cm} & \linead{4.2cm}\\
5. Intergeneracional & A mi abuela/abuelo le diría \linead{6.5cm} & \linead{4.2cm}\\
\end{tabularx}
\normalsize

\begin{infoband}{milcVino}
\textbf{Para la próxima sustentación.} Vuelve a esta tabla en seis meses. Verás que las respuestas cambian — no porque lo aprendido cambie, sino porque tú cambias. La evaluación liberadora es una conversación contigo a través del tiempo.
\end{infoband}


\puente{Cerramos con 3 voces que te ayudan a entender por qué conocer las piezas te da poder.}

\begin{titlebox}{milcGradeDark}
Triángulo de pensamiento · cierre filosófico
\end{titlebox}

\begin{softbox}{milcVino}{milcGris}{Dussel · lente del nosotros}
\textit{````Conocer el nombre de algo es el primer paso para perderle el miedo.''''}

Antes de hoy, el interior del computador era un misterio. \textbf{Ahora ya conoces el nombre de las 6 piezas más importantes}. Eso no es teoría: es poder. La próxima vez que un técnico te diga \emph{``el problema es la fuente de poder''}, ya no te suena chino: sabes qué es y por qué importa. \textbf{Los nombres son llaves}: cada nombre que aprendes te abre una puerta.

\textbf{¿De cuántas cosas dependo todos los días sin saber su nombre? ¿Cuál sería la próxima que querría aprender a nombrar?}
\end{softbox}
\linead{\linewidth}\\[1mm]
\linead{\linewidth}

\begin{softbox}{milcMostaza}{milcGris}{Estoicismo · Epicteto · cuidado interior}
\textit{````Lo que se entiende, se domina. Lo que no se entiende, te domina a ti.''''}

Cuando no entiendes cómo funciona algo, ese algo te genera ansiedad: \emph{``¿y si se daña?''}, \emph{``¿qué hago si no funciona?''}. Cuando lo entiendes, te calmas. \textbf{No es magia: es comprensión}. Saber qué hace cada pieza del gabinete te quita esa ansiedad. Si algo falla, ahora tienes vocabulario para preguntar y mapa para buscar.

\textbf{¿Qué cosas me asustan solo porque no las entiendo? ¿Cuál vale la pena empezar a entender hoy?}
\end{softbox}
\linead{\linewidth}\\[1mm]
\linead{\linewidth}

\begin{softbox}{milcTurquesa}{milcGris}{Floridi · lente de la infoesfera}
\textit{````Saber el nombre de las cosas técnicas es la primera forma de ciudadanía en una sociedad de máquinas.''''}

Vivimos rodeados de máquinas. Pero \textbf{la mayoría de las personas no sabe cómo funcionan por dentro}. Eso las hace dependientes: dependen del técnico, del youtuber, del vendedor que les dice qué comprar. Cuando tú aprendes los nombres y las funciones, dejas de depender. Eres un \emph{ciudadano técnicamente alfabetizado}: una de las personas que entienden el mundo donde viven.

\textbf{¿Quiero ser dependiente de los expertos para todo, o quiero entender lo suficiente para decidir por mí?}
\end{softbox}
\linead{\linewidth}\\[1mm]
\linead{\linewidth}

\begin{titlebox}{milcVino}
Compromiso final
\end{titlebox}

\small
\begin{tabularx}{\textwidth}{>{\RaggedRight\arraybackslash}p{3.0cm}YYY}
\rowcolor{milcVino}\color{white}\bfseries Criterio & \color{white}\bfseries Lo logré & \color{white}\bfseries En proceso & \color{white}\bfseries Necesito apoyo\\
Comprensión del tema & Explico ideas centrales con mis palabras. & Reconozco ideas pero falta precisión. & Necesito repasar conceptos base.\\
Pensamiento computacional & Apliqué los 4 pilares al diseñar mi producto. & Apliqué algunos pilares. & Necesito releer la fase 2.\\
Aplicación y producto & Mi producto cumple los criterios y se puede revisar. & Producto funcional con ajustes pendientes. & Aún en construcción.\\
Reflexión 5 dimensiones & Respondí cada dimensión con honestidad. & Respondí algunas con detalle. & Mis respuestas son muy generales.\\
Triángulo de pensamiento & Conecté las 3 voces con mi trabajo real. & Conecté al menos una. & No relaciono las voces con mi proceso.\\
\end{tabularx}
\normalsize

\textbf{Hoy entendí que:}\\[1mm]
\linead{\linewidth}\\[1mm]
\linead{\linewidth}

\textbf{Mi compromiso para la próxima sesión es:}\\[1mm]
\linead{\linewidth}\\[1mm]
\linead{\linewidth}

\begin{softbox}{milcMostaza}{milcGris}{Cita de despedida · Marco Aurelio}
\textit{``El obstáculo en el camino se vuelve el camino. Lo que estorba al camino se vuelve camino.''}\\[1mm]
Cada error de hoy, bien observado, se convierte en pieza del camino que pisarás mañana.
\end{softbox}

\small
\begin{tabularx}{\textwidth}{>{\bfseries}p{3.6cm}Y}
\rowcolor{milcGradeDark!12}Nombre completo: & \linead{10cm}\\
Fecha: & \linead{4cm} \quad Curso/grupo: \linead{4cm}\\
Mi firma: & \linead{9cm}\\
\end{tabularx}
\normalsize

\begin{infoband}{milcGradeDark}
\textbf{Para la próxima sesión.} Esta semana, si entras a la sala de sistemas, observa los gabinetes desde afuera: \textbf{¿qué cables ves?} ¿Cuáles van a la fuente de poder? ¿Cuáles van al monitor (eso es salida)? La próxima clase trabajamos los \textbf{periféricos}: las piezas que están \emph{afuera} del gabinete --- teclado, ratón, pantalla, cámara, impresora. Hoy viste lo interno; la próxima ves lo externo.
\end{infoband}

\end{document}
